\author{OTS}
\documentclass[fontsize=16pt]{mybib}
\renewcommand{\DFA}{12mm}
\renewcommand{\DZA}{1.3}

\graphicspath{{../../../assets/images/}}
\setheaderlogo{mnr-green.png}

% Absatzformatierung
\setlength{\parindent}{0pt}
\setlength{\parskip}{0.6em}
\renewcommand{\arraystretch}{0.8}



\title{Nehemia 1}
\author{Lothar Schmid}
\date{}

\begin{document}
\mytitlepage
\tableofcontents
\flushbottom
\newpage

\section{Teil 0: Einleitung}
Bevor wir uns die Taten von Nehemia näher anschauen, wollen wir uns ein Übersicht über das Buch Nehemia verschaffen. Auch wenn ein Teil des Buches in der ICH Form geschreiben wurde, wird angenommen, dass dieses Buch von Esra geschrieben wurde und Esra auf Aufzeichnungen von Nehemia zurückgegriffen hat. Die Bücher Esra und Nehemia sind in der Thora ein Buch. Dass aber in beiden Bücher ein praktisch identisches Geschlechtsregister aufgeführt wird, zeigt aber eher, dass es sich zu Beginn wohl um zwei Bücher gehandelt hat. Macht ja keinen Sinn, zwei identische Namenslisten aufzuführen.

Der Zeitraum in dem sich diese Handlungen ablaufen kann man von 538 - 430 \vChr eingegrenzt werden. Kurz nachdem die Persier das Land von den Babiloniern erobert hatte, fand auch die erste Rückführung von Juden nach Israel statt. In 
\begin{bibelbox}{Jer}{25:11-12}
    \hh{11}und dieses ganze Land soll zu Trümmerhaufen, zur Wüste werden, und diese Völker sollen dem König von Babel dienen, 70 Jahre lang. \hh{12}Und es wird geschehen, wenn die 70 Jahre vollendet sind, dann will ich an dem König von Babel und an jenem Volk ihre Schuld heimsuchen, spricht der HERR, auch am Land der Chaldäer, und ich will es zur ewigen Wüste machen.
\end{bibelbox}
\begin{bibelbox}{Jer}{29:10}
    \hh{10}Fürwahr, so spricht der \herr : Wenn die 70 Jahre für Babel gänzlich erfüllt sind, werde ich mich euer annehmen und mein gutes Wort, euch an diesen Ort zurückzubringen, an euch erfüllen.
\end{bibelbox}
Dies verkündete Jeremia im Auftrag Gottes kurz bevor die Juden von den Babiloniern verschleppt wurde. Intressant in dem Zusammenhang ist auch, dass Gott den Juden den Rat gegeben hat, in dem fremden Land sich zu integrieren Familien zu gründen und den Frieden zu suchen. Sagt nicht auch Paulus, dass wir Fremde auf dieser Welt sind? 
\begin{bibelbox}{Phil}{3:20}
    \hh{20}Unser Bürgerrecht aber ist im Himmel, von woher wir auch den Herrn Jesus Christus erwarten als den Retter, 
\end{bibelbox}
Und Petrus, dass wir uns benehmen sollen?
\begin{bibelbox}{IPetr}{2:11-12}
    \hh{11}Geliebte, ich ermahne euch als Gäste und Fremdlinge[4]: Enthaltet euch der fleischlichen Begierden, die gegen die Seele streiten; 
    \hh{12} und führt einen guten Wandel unter den Heiden, damit sie da, wo sie euch als Übeltäter verleumden, doch aufgrund der guten Werke, die sie gesehen haben, Gott preisen am Tag der Untersuchung.
\end{bibelbox}
\end{document}
